% Options for packages loaded elsewhere
\PassOptionsToPackage{unicode}{hyperref}
\PassOptionsToPackage{hyphens}{url}
\PassOptionsToPackage{dvipsnames,svgnames,x11names}{xcolor}
%
\documentclass[
  12pt]{article}

\usepackage{amsmath,amssymb}
\usepackage{iftex}
\ifPDFTeX
  \usepackage[T1]{fontenc}
  \usepackage[utf8]{inputenc}
  \usepackage{textcomp} % provide euro and other symbols
\else % if luatex or xetex
  \usepackage{unicode-math}
  \defaultfontfeatures{Scale=MatchLowercase}
  \defaultfontfeatures[\rmfamily]{Ligatures=TeX,Scale=1}
\fi
\usepackage{lmodern}
\ifPDFTeX\else  
    % xetex/luatex font selection
\fi
% Use upquote if available, for straight quotes in verbatim environments
\IfFileExists{upquote.sty}{\usepackage{upquote}}{}
\IfFileExists{microtype.sty}{% use microtype if available
  \usepackage[]{microtype}
  \UseMicrotypeSet[protrusion]{basicmath} % disable protrusion for tt fonts
}{}
\makeatletter
\@ifundefined{KOMAClassName}{% if non-KOMA class
  \IfFileExists{parskip.sty}{%
    \usepackage{parskip}
  }{% else
    \setlength{\parindent}{0pt}
    \setlength{\parskip}{6pt plus 2pt minus 1pt}}
}{% if KOMA class
  \KOMAoptions{parskip=half}}
\makeatother
\usepackage{xcolor}
\setlength{\emergencystretch}{3em} % prevent overfull lines
\setcounter{secnumdepth}{5}
% Make \paragraph and \subparagraph free-standing
\makeatletter
\ifx\paragraph\undefined\else
  \let\oldparagraph\paragraph
  \renewcommand{\paragraph}{
    \@ifstar
      \xxxParagraphStar
      \xxxParagraphNoStar
  }
  \newcommand{\xxxParagraphStar}[1]{\oldparagraph*{#1}\mbox{}}
  \newcommand{\xxxParagraphNoStar}[1]{\oldparagraph{#1}\mbox{}}
\fi
\ifx\subparagraph\undefined\else
  \let\oldsubparagraph\subparagraph
  \renewcommand{\subparagraph}{
    \@ifstar
      \xxxSubParagraphStar
      \xxxSubParagraphNoStar
  }
  \newcommand{\xxxSubParagraphStar}[1]{\oldsubparagraph*{#1}\mbox{}}
  \newcommand{\xxxSubParagraphNoStar}[1]{\oldsubparagraph{#1}\mbox{}}
\fi
\makeatother


\providecommand{\tightlist}{%
  \setlength{\itemsep}{0pt}\setlength{\parskip}{0pt}}\usepackage{longtable,booktabs,array}
\usepackage{calc} % for calculating minipage widths
% Correct order of tables after \paragraph or \subparagraph
\usepackage{etoolbox}
\makeatletter
\patchcmd\longtable{\par}{\if@noskipsec\mbox{}\fi\par}{}{}
\makeatother
% Allow footnotes in longtable head/foot
\IfFileExists{footnotehyper.sty}{\usepackage{footnotehyper}}{\usepackage{footnote}}
\makesavenoteenv{longtable}
\usepackage{graphicx}
\makeatletter
\def\maxwidth{\ifdim\Gin@nat@width>\linewidth\linewidth\else\Gin@nat@width\fi}
\def\maxheight{\ifdim\Gin@nat@height>\textheight\textheight\else\Gin@nat@height\fi}
\makeatother
% Scale images if necessary, so that they will not overflow the page
% margins by default, and it is still possible to overwrite the defaults
% using explicit options in \includegraphics[width, height, ...]{}
\setkeys{Gin}{width=\maxwidth,height=\maxheight,keepaspectratio}
% Set default figure placement to htbp
\makeatletter
\def\fps@figure{htbp}
\makeatother

\addtolength{\oddsidemargin}{-.5in}%
\addtolength{\evensidemargin}{-1in}%
\addtolength{\textwidth}{1in}%
\addtolength{\textheight}{1.7in}%
\addtolength{\topmargin}{-1in}%
\makeatletter
\@ifpackageloaded{caption}{}{\usepackage{caption}}
\AtBeginDocument{%
\ifdefined\contentsname
  \renewcommand*\contentsname{Table of contents}
\else
  \newcommand\contentsname{Table of contents}
\fi
\ifdefined\listfigurename
  \renewcommand*\listfigurename{List of Figures}
\else
  \newcommand\listfigurename{List of Figures}
\fi
\ifdefined\listtablename
  \renewcommand*\listtablename{List of Tables}
\else
  \newcommand\listtablename{List of Tables}
\fi
\ifdefined\figurename
  \renewcommand*\figurename{Figure}
\else
  \newcommand\figurename{Figure}
\fi
\ifdefined\tablename
  \renewcommand*\tablename{Table}
\else
  \newcommand\tablename{Table}
\fi
}
\@ifpackageloaded{float}{}{\usepackage{float}}
\floatstyle{ruled}
\@ifundefined{c@chapter}{\newfloat{codelisting}{h}{lop}}{\newfloat{codelisting}{h}{lop}[chapter]}
\floatname{codelisting}{Listing}
\newcommand*\listoflistings{\listof{codelisting}{List of Listings}}
\makeatother
\makeatletter
\makeatother
\makeatletter
\@ifpackageloaded{caption}{}{\usepackage{caption}}
\@ifpackageloaded{subcaption}{}{\usepackage{subcaption}}
\makeatother

\ifLuaTeX
  \usepackage{selnolig}  % disable illegal ligatures
\fi
\usepackage[]{natbib}
\bibliographystyle{agsm}
\usepackage{bookmark}

\IfFileExists{xurl.sty}{\usepackage{xurl}}{} % add URL line breaks if available
\urlstyle{same} % disable monospaced font for URLs
\hypersetup{
  pdftitle={Introduction},
  pdfauthor={Kevin Linares},
  pdfkeywords={latent class analysis, complex survey
data, pseudo-maximum likelihood, replicate variance
estimation, three-step estimation, institutional trust},
  colorlinks=true,
  linkcolor={blue},
  filecolor={Maroon},
  citecolor={Blue},
  urlcolor={Blue},
  pdfcreator={LaTeX via pandoc}}



\begin{document}


\def\spacingset#1{\renewcommand{\baselinestretch}%
{#1}\small\normalsize} \spacingset{1}


%%%%%%%%%%%%%%%%%%%%%%%%%%%%%%%%%%%%%%%%%%%%%%%%%%%%%%%%%%%%%%%%%%%%%%%%%%%%%%

\date{July 20, 2026}
\title{\bf Introduction}
\author{
Kevin Linares\\
Joint Program in Survey Methodology, University of Maryland\\
}
\maketitle

\bigskip
\bigskip
\begin{abstract}
Latent class analysis (LCA) of categorical items collected under a
stratified, clustered, weighted sample design requires three things that
no single R package currently provides together: a design-weighted
objective so that estimates describe the population rather than the
achieved sample, model selection statistics placed on a scale consistent
with that objective, and variance estimation that respects the design in
the absence of a usable information matrix. We present an integrated
pipeline meeting all three requirements. Parameters are estimated by
pseudo-maximum likelihood through a weighted EM algorithm; information
criteria are calibrated by weight normalization; and all uncertainty
statements derive from stratified delete-one jackknife (JKn)
replication, with replicate refits warm-started and aligned to the
full-sample solution. Segment composition across external variables is
estimated by the BCH three-step correction with design-based replicate
intervals, and the unweighted special case is validated to exact
agreement against an independent maximum likelihood implementation. We
illustrate the methodology with thirteen institutional trust items from
the 2023 AmericasBarometer survey of Mexico, identifying five trust
segments whose composition varies sharply with presidential approval and
satisfaction with democracy. Class naming is drafted by a large language
model operating under a certified, frozen prompt instrument with full
analyst override.
\end{abstract}

\noindent%
{\it Keywords:} latent class analysis, complex survey
data, pseudo-maximum likelihood, replicate variance
estimation, three-step estimation, institutional trust
\vfill

\newpage
\spacingset{1.9} % DON'T change the spacing!


Latent class analysis identifies unobserved population subgroups from
response patterns over categorical indicators
\citep{lazarsfeld1968, goodman1974}. When the indicators come from a
complex probability sample, three methodological requirements arise
jointly. First, the likelihood must be weighted so that the recovered
classes and their prevalences describe the target population rather than
the achieved sample \citep{patterson2002, asparouhov2005}. Second, model
selection statistics must sit on a scale consistent with the weighted
objective, since base weights rarely sum to the sample size and the
pseudo-log-likelihood is otherwise miscalibrated against the penalty
term \citep{lumley2015}. Third, variance estimation must respect the
stratification and clustering of the design; under a pseudo-likelihood
the inverse information matrix is not a valid variance estimator, so
design-based replication is the natural device
\citep{wolter2007, rust1996}.

No R package meets these three requirements together with a path from
measurement model to labelled analytic data product. The \texttt{poLCA}
package \citep{linzer2011} fits unweighted maximum likelihood only. The
\texttt{survey} package \citep{lumley2004} supplies design declaration
and replication machinery but no mixture models. The closest existing
software, the \texttt{baysc} package \citep{wu2024, wu2026joss}, fits
the same weighted objective within a Bayesian pseudo-posterior with a
sandwich-based variance adjustment; the pipeline described here is its
frequentist counterpart, carrying replicate-based variance through model
selection, measurement-model inference, bias-adjusted profiling on
external variables, prediction for item-partial respondents, and a
labelled data file for downstream analysts. Commercial options exist,
notably Mplus with \texttt{TYPE\ =\ MIXTURE\ COMPLEX} and Latent GOLD,
and Stata's \texttt{gsem} under \texttt{svy}; the contribution here is
an open, auditable implementation in R in which every estimator
corresponds line for line to documented source code, with the unweighted
special case validated against an independent implementation.

Throughout, we refer to latent classes as segments in all user-facing
output, a terminological preference of applied audiences; the
statistical object is unchanged.

We illustrate the methodology with an analysis of institutional trust in
Mexico using the 2023 AmericasBarometer \citep{lapop2023}. Institutional
trust batteries are a natural LCA application: respondents rarely trust
or distrust all institutions uniformly, and the substantive interest
lies in the configuration of trust across institutions, which a
unidimensional scale conceals. Five segments emerge, ranging from
generalized distrust to broad confidence, with two intermediate profiles
distinguished chiefly by their treatment of the president and the armed
forces relative to representative and justice institutions.

Data

The 2023 AmericasBarometer round in Mexico interviewed 1,622 voting-age
adults under a stratified, clustered design with 4 strata and 130
primary sampling units, released with base weights \citep{lapop2023}. We
analyze the thirteen-item institutional trust battery: trust in the
justice system, the electoral tribunal, the armed forces, the
legislature, the public ministry, the national police, the federal
auditor, political parties, the president, the supreme court, the
respondent's municipality, the media, and elections. Each item is
answered on a seven-point scale anchored at 1, none at all
(\emph{nada}), and 7, a lot (\emph{mucho}); intermediate points are
unlabelled. Nonresponse codes are recoded to missing during preparation,
and each item is recoded to consecutive integers over its substantive
categories.

Estimation uses complete cases on the thirteen items and the profiling
covariates: 1,430 of 1,622 respondents (88.2 percent) enter the analysis
sample. This is a deliberate design choice rather than an oversight;
Section~\ref{sec-assignment} describes how the fitted model then assigns
segment membership to item-partial respondents in the preserved full
file, subject to a minimum-information floor, so the complete-case
restriction affects estimation but not the coverage of the delivered
data product. Profiling covariates are age category, sex, education,
urban or rural residence, employment status, satisfaction with
democracy, and evaluation of the president's job performance.

Methods

\section{Model}\label{sec-model}

Let \(\mathbf{y}_i = (y_{i1}, \dots, y_{iJ})\) denote respondent \(i\)'s
responses to \(J\) nominal items, item \(j\) having \(C_j\) categories.
With \(K\) latent segments, prevalences \(\pi_k\), and conditional
response probabilities \(\rho_{jck} = P(y_{ij} = c \mid k)\), the
marginal probability of a response pattern is

\begin{equation}\phantomsection\label{eq-lca}{
P(\mathbf{y}_i) \;=\; \sum_{k=1}^{K} \pi_k \prod_{j=1}^{J} \rho_{j\,y_{ij}\,k},
\qquad \sum_{c=1}^{C_j} \rho_{jck} = 1 .
}\end{equation}

Conditional independence of the items given segment membership, local
independence, is the load-bearing assumption;
Section~\ref{sec-diagnostics} probes where it strains and
Section~\ref{sec-assignment} exploits it for prediction under item
nonresponse.

\section{Design-weighted pseudo-maximum likelihood}\label{sec-pml}

Under a complex design with base weights \(w_i\), estimation maximizes
the pseudo-log-likelihood

\begin{equation}\phantomsection\label{eq-pll}{
\ell_w(\boldsymbol\theta) \;=\; \sum_{i=1}^{n} w_i \,
\log\!\Big( \sum_{k=1}^{K} \pi_k \prod_{j=1}^{J} \rho_{j\,y_{ij}\,k} \Big),
}\end{equation}

the standard device for LCA under complex sampling
\citep{patterson2002, skinner1989}. The EM algorithm
\citep{dempster1977} alternates an E-step computing posterior membership
probabilities with the log-sum-exp device for numerical stability,

\begin{equation}\phantomsection\label{eq-estep}{
p_{ik} \;=\; \frac{\pi_k \prod_j \rho_{j\,y_{ij}\,k}}
{\sum_{l} \pi_l \prod_j \rho_{j\,y_{ij}\,l}},
}\end{equation}

and an M-step in which the weights enter every sufficient statistic,

\begin{equation}\phantomsection\label{eq-mstep}{
\hat\pi_k = \frac{\sum_i w_i\, p_{ik}}{\sum_i w_i}, \qquad
\hat\rho_{jck} = \frac{\sum_i w_i\, p_{ik}\, \mathbf{1}(y_{ij} = c)}
{\sum_i w_i\, p_{ik}} .
}\end{equation}

Iteration stops when \(|\Delta \ell_w| < \varepsilon\,(|\ell_w| + 1)\).
During enumeration \(\varepsilon = 10^{-8}\) with at most 800 iterations
per start; the chosen model is refit at terminal precision,
\(\varepsilon = 10^{-10}\) with at most 5,000 iterations, matching the
stopping behavior of the reference implementation used for validation so
that the comparison in Section~\ref{sec-validation} contrasts optima
rather than stopping rules. Because the pseudo-likelihood of a finite
mixture can be multimodal with near-tied optima, each candidate model is
fit from twenty random starts with deterministic seeding, and the
reported solution is the best mode found; this claim discipline is
stated rather than hidden. Point estimates are invariant to the scale of
the weights.

\section{Model selection}\label{sec-selection}

Information criteria compare the maximized log-likelihood against a
sample-size penalty, and the two must sit on the same scale. Base
weights rarely sum to \(n\); because \(\ell_w\) is linear in the
weights, multiplying it by \(n / \sum_i w_i\) is exactly equivalent to
fitting with weights rescaled to sum to \(n\), leaving point estimates
untouched and calibrating only model selection {[}\citet{lumley2015}
develop the general design-based treatment; the sum-to-\(n\) convention
matches the default weight normalization in Mplus, making criteria
comparable across software{]}. With
\(\ell^{*}_w = \ell_w \cdot n / \sum_i w_i\) and model dimension

\begin{equation}\phantomsection\label{eq-df}{
d_K = (K - 1) + K \sum_{j=1}^{J} (C_j - 1),
}\end{equation}

we report \(\mathrm{BIC} = -2\ell^{*}_w + d_K \log n\) and
\(\mathrm{AIC} = -2\ell^{*}_w + 2 d_K\), alongside the relative entropy

\begin{equation}\phantomsection\label{eq-entropy}{
R^2_E \;=\; 1 - \frac{-\sum_{i,k} p_{ik} \log p_{ik}}{\,n \log K\,},
}\end{equation}

which indexes classification sharpness \citep{ramaswamy1993}. The
bootstrapped likelihood ratio test, often recommended for class
enumeration under simple random sampling \citep{nylund2007}, is
deliberately omitted: its parametric bootstrap presumes independent
observations, which fails under a design-weighted pseudo-likelihood on a
clustered sample, matching the equivalent restriction imposed by Mplus
under complex-design analysis. The number of segments is never selected
automatically. The analyst chooses \(K\) by iterating over these
statistics and the diagnostics of Section~\ref{sec-diagnostics}, and
fixes it in configuration; Section~\ref{sec-enumeration} reports that
reasoning for the present application.

\section{Replicate variance estimation}\label{sec-variance}

A pseudo-likelihood provides no usable information matrix, so all
uncertainty statements come from stratified delete-one-PSU jackknife
(JKn) replication \citep{wolter2007, rust1996, lumley2004}. With \(n_h\)
PSUs in stratum \(h\), deleting PSU \(j\) of stratum \(h\) and
reweighting its stratum-mates by \(n_h/(n_h - 1)\) yields the replicate
estimate \(\hat\theta_{(hj)}\), and

\begin{equation}\phantomsection\label{eq-jkn}{
\widehat{V}(\hat\theta) \;=\; \sum_{h} \frac{n_h - 1}{n_h}
\sum_{j=1}^{n_h} \big( \hat\theta_{(hj)} - \hat\theta \big)^2 .
}\end{equation}

Each replicate refits the weighted EM on the replicate weights,
warm-started from the full-sample solution, so the resulting variance is
within-mode by construction: it carries design variance, not uncertainty
about mode selection. Label switching across replicates is resolved by
Hungarian assignment \citep{kuhn1955, hornik2005} on the conditional
response profiles, aligning every replicate to the full-sample fit
before differencing.

Confidence intervals for every probability, the prevalences \(\pi_k\)
and the conditional probabilities \(\rho_{jck}\), are formed on the
logit scale with the delta-method standard error
\(\mathrm{se}_{\mathrm{logit}} = \mathrm{se} / \{\hat p (1 - \hat p)\}\)
and back-transformed, which keeps them inside \((0, 1)\) without
truncation and improves finite-sample coverage away from one half. The
critical value is \(t_{0.975, \nu}\) with \(\nu\) the design degrees of
freedom, the number of PSUs minus the number of strata, here
\(130 - 4 = 126\), rather than the normal 1.96. Cells estimated within
\(10^{-3}\) of the boundary have no meaningful logit-scale standard
error and instead receive Wald intervals truncated to \([0, 1]\),
flagged as such in the delivered tables. The same interval construction
serves the domain estimates of Section~\ref{sec-bch}, so a single method
runs through the entire pipeline.

\section{Model-fit diagnostics}\label{sec-diagnostics}

Item discrimination summarizes how much item \(j\) separates the
segments as the mean, over segment pairs, of the total variation
distance between their conditional response distributions:

\begin{equation}\phantomsection\label{eq-discrim}{
D_j \;=\; \frac{2}{K(K-1)} \sum_{k < l} \frac{1}{2} \sum_{c=1}^{C_j}
\big| \rho_{jck} - \rho_{jcl} \big| .
}\end{equation}

\(D_j = 0\) means every segment answers item \(j\) identically, a
removal candidate for a sensitivity refit; values near 1 indicate sharp
separation.

Discrimination is an effect size; a companion screen asks whether each
item earns its class-dependent parameters by BIC. For item \(j\), the
fitted model is compared on the same respondents against a restricted
model in which item \(j\) is made class-independent, its response
distribution shared across segments. Under local independence the
restricted model factorizes into an LCA on the remaining items times
item \(j\)'s weighted marginal, so its pseudo-log-likelihood is
available without a constrained refit as
\(\ell_{-j} + \ell_j^{\mathrm{marg}}\), its dimension is
\(d_K^{(-j)} + (C_j - 1)\), and its BIC on the sum-to-\(n\) scale is
directly comparable to the full model's. We report
\(\Delta\mathrm{BIC}_j = \mathrm{BIC}_{\mathrm{restricted}} - \mathrm{BIC}_{\mathrm{full}}\);
a negative value means BIC prefers dropping item \(j\)'s class
dependence and flags the item.

Local independence is probed pairwise with a bivariate residual (BVR).
For items \(a, b\), let \(\widehat{P}_{ab}(c,d)\) be the design-weighted
observed proportion of the joint response \((c,d)\) and
\(\Pi_{ab} = \boldsymbol\rho_a \operatorname{diag}(\boldsymbol\pi)\, \boldsymbol\rho_b^{\top}\)
the model-implied joint distribution. The statistic

\begin{equation}\phantomsection\label{eq-bvr}{
\mathrm{BVR}_{ab} \;=\; \frac{1}{(C_a - 1)(C_b - 1)}
\; n \sum_{c,d} \frac{\big( \widehat{P}_{ab}(c,d) - \Pi_{ab}(c,d) \big)^2}{\Pi_{ab}(c,d)}
}\end{equation}

is a Pearson-type residual per degree of freedom. In unweighted maximum
likelihood settings its null distribution is already not chi-square and
bootstrap calibration is advised \citep{oberski2013}; under weighting
and pseudo-likelihood no reference distribution is valid, so we report
BVR strictly as a descriptive ranking of the item pairs the model
explains least well, with large persistent values arguing for item
removal, merging, or \(K+1\).

\section{Validation against an independent
implementation}\label{sec-validation}

At unit weights, Equation~\ref{eq-pll} reduces to the ordinary
log-likelihood, so the weighted EM and \texttt{poLCA} \citep{linzer2011}
maximize the same function and must agree up to optimizer tolerance and
mode selection. The pipeline enforces this as a runtime check. The
unit-weight EM solution is obtained from the same multi-start protocol
plus one additional candidate warm-started at the \texttt{poLCA}
solution itself; since both maximize the identical function, whichever
basin attains the higher log-likelihood is retained. The retained
solution is aligned to the \texttt{poLCA} solution by Hungarian matching
on response profiles, and maximum absolute gaps in prevalences and
conditional probabilities are computed and frozen to disk with a single
pass or fail verdict at a 0.005 threshold on both gaps, printed
identically in every downstream document. Where a gap appears,
comparison of the two log-likelihoods separates flat-ridge drift under
the stopping rule, an inferior \texttt{poLCA} mode, and an inferior EM
mode calling for more starts. The same reasoning positions the pipeline
against Mplus (\texttt{TYPE\ =\ MIXTURE\ COMPLEX}) and Stata
(\texttt{gsem} under \texttt{svy}): identical objective and weight
convention imply matching point estimates and criteria up to mode
selection, with linearization rather than replication standard errors,
both design-consistent.

\section{Segment assignment under partial
response}\label{sec-assignment}

Membership is assigned from the posterior evaluated at the fitted
parameters. Local independence makes assignment honest for item-partial
respondents: a missing item drops out of the within-segment product, so
for respondent \(i\) with answered item set \(A_i\),

\begin{equation}\phantomsection\label{eq-partial}{
p_{ik}^{\mathrm{obs}} \;=\;
\frac{\pi_k \prod_{j \in A_i} \rho_{j\,y_{ij}\,k}}
{\sum_l \pi_l \prod_{j \in A_i} \rho_{j\,y_{ij}\,l}},
}\end{equation}

and the assigned segment is \(\arg\max_k p_{ik}^{\mathrm{obs}}\). An
information floor \(m^{*}\), configured at seven of the thirteen items
here, is enforced: respondents answering fewer than \(m^{*}\) items
receive no assignment rather than a weakly informed one. The model is
fit on complete cases only; assignment extends to the preserved full
frame, and the exported labelled data record which respondents were
assigned.

\section{Segment composition on external variables: the BCH
correction}\label{sec-bch}

Cross-classifying modal assignments against external variables
understates segment differences because assignment error, even at high
entropy, attenuates composition toward the marginal \citep{bolck2004}.
We therefore estimate segment prevalence within demographic domains by
the BCH three-step correction \citep{bolck2004, vermunt2010}, adapted to
the design-weighted setting. Step one is the fitted measurement model.
Step two computes the modal assignment \(W_i = \arg\max_k p_{ik}\) and
the design-weighted classification-error matrix \(\mathbf{D}\) with
entries

\begin{equation}\phantomsection\label{eq-bchD}{
D_{ks} \;=\; \widehat{P}(W = s \mid X = k) \;=\;
\frac{\sum_i w_i\, p_{ik}\, \mathbf{1}(W_i = s)}{\sum_i w_i\, p_{ik}},
}\end{equation}

the weighted posterior mass of true segment \(k\) landing in assigned
segment \(s\); rows of \(\mathbf{D}\) sum to one. Step three replaces
each respondent's hard assignment with the correction vector
\(\mathbf{u}_i = \mathbf{e}_{W_i}^{\top} \mathbf{D}^{-1}\), row \(W_i\)
of \(\mathbf{D}^{-1}\). Weighted by \(w_i\), the \(\mathbf{u}_i\) are
unbiased for true-segment membership, so the BCH-corrected prevalence of
segment \(k\) within domain \(d\) is the design-based domain mean

\begin{equation}\phantomsection\label{eq-bchdom}{
\hat p_{dk} \;=\; \frac{\sum_{i \in d} w_i\, u_{ik}}{\sum_{i \in d} w_i}.
}\end{equation}

Two properties of the correction deserve note. Entries of
\(\mathbf{u}_i\) can be negative, so individual cell estimates can fall
slightly outside \([0, 1]\); this is a documented property of the
method, not an error. And each \(\mathbf{u}_i\) sums to one because the
rows of \(\mathbf{D}\) do, so within every domain the estimated
prevalences still sum to one across segments.

Variance is estimated with the same JKn replicates: the posterior matrix
and modal assignments are held fixed at the full-sample measurement
model, while \(\mathbf{D}\) and the correction weights are rebuilt
inside every replicate from the replicate weights, so the intervals
carry the design variance of the classification table and the
cross-tabulation jointly. Holding step one fixed is the standard
treatment; the resulting understatement of variance is minor when
classification is sharp \citep{bakk2014, bakk2013}, and the present
application, with relative entropy 0.907, is squarely in that favorable
regime. Intervals follow Section~\ref{sec-variance}, logit-scale with
the design \(t\) critical value; cells outside \((0, 1)\), where the
logit transform is undefined, fall back to truncated Wald intervals and
are footnoted in the delivered tables.

\section{LLM-assisted, analyst-controlled labeling}\label{sec-labeling}

Segment naming is a bottleneck in applied LCA reporting, and recent
evidence suggests large language models perform annotation-adjacent
tasks reliably when tightly instructed \citep{gilardi2023, ziems2024}.
The pipeline drafts segment names with a language model under a fixed
instrument: a persona and rule set instructing interpretation strictly
from the displayed conditional response probabilities and item wording,
one isolated call per segment (joint prompts empirically confused
near-neighbor segments), an optional survey-context paragraph fenced to
referent resolution only, and a collision-gated harmonization pass that
may edit labels, never descriptions. Labels follow one rule: a label
file on disk is used when present and validated against \(K\); otherwise
the model drafts once and the result is frozen to that file, which the
analyst edits to take over naming. The freeze file, not a sampling seed,
is the reproducibility mechanism.

The instrument is treated as a measurement device. A certification
harness plants a known-meaning model, including near-neighbor segments,
a deliberately diffuse segment, nonresponse-labelled items, and a
fictional survey context, and checks output format, grounding in the
displayed probabilities, hedging on diffuse profiles, persona integrity,
context-leak fencing, harmonizer properties, and adoption of the segment
terminology; any edit to the instrument obliges recertification. Labels
are drafts, verified by the analyst against the conditional response
profiles before use, and the delivered tables carry the label file's
contents, not live model output.

Results

\section{Enumeration and the choice of K}\label{sec-enumeration}

Table~\ref{tbl-enum} reports the weighted fit statistics for \(K = 2\)
through \(12\), each fit from twenty seeded random starts; all
candidates converged. Figure~\ref{fig-enum} plots weighted BIC and
relative entropy.

\begin{longtable}[]{@{}rrrrrr@{}}

\caption{\label{tbl-enum}Weighted fit statistics by number of segments.
The log-likelihood is on the sum-to-n scale; entropy is the relative
entropy of the posterior classification.}

\tabularnewline

\toprule\noalign{}
K & Log-likelihood & df & AIC & BIC & Entropy \\
\midrule\noalign{}
\endhead
\bottomrule\noalign{}
\endlastfoot
2 & -31579.7 & 157 & 63473.3 & 64300.0 & 0.904 \\
3 & -30176.2 & 236 & 60824.3 & 62067.0 & 0.913 \\
4 & -29447.7 & 315 & 59525.5 & 61184.1 & 0.921 \\
5 & -29047.7 & 394 & 58883.4 & 60957.9 & 0.907 \\
6 & -28751.8 & 473 & 58449.5 & 60940.1 & 0.895 \\
7 & -28537.7 & 552 & 58179.3 & 61085.8 & 0.893 \\
8 & -28390.2 & 631 & 58042.5 & 61365.0 & 0.890 \\
9 & -28242.7 & 710 & 57905.4 & 61643.8 & 0.890 \\
10 & -28123.2 & 789 & 57824.4 & 61978.8 & 0.894 \\
11 & -28024.0 & 868 & 57784.0 & 62354.4 & 0.902 \\
12 & -27927.8 & 947 & 57749.6 & 62736.0 & 0.899 \\

\end{longtable}

\begin{figure}

\centering{

\includegraphics{survey_lca_paper_files/figure-pdf/fig-enum-1.pdf}

}

\caption{\label{fig-enum}Enumeration diagnostics across K: weighted BIC
and relative entropy. BIC is essentially flat between K = 5 and K = 6;
entropy peaks at K = 4 and declines thereafter.}

\end{figure}%

Weighted BIC drops steeply through \(K = 5\) (60,957.9), reaches its
minimum at \(K = 6\) (60,940.1), and rises monotonically thereafter. The
difference between the two, 17.8, is small against the preceding drops
of 226.2 (\(4 \to 5\)) and 882.9 (\(3 \to 4\)): the criterion is
effectively flat over the \(\{5, 6\}\) region. AIC continues to decline
through the whole range, its usual overextraction under large \(n\)
\citep{nylund2007}. Relative entropy peaks at \(K = 4\) (0.921) and
remains high at \(K = 5\) (0.907); no larger \(K\) exceeds 0.91. Within
the flat BIC region we fixed \(K = 5\) on parsimony, classification
sharpness, and interpretability: the five-segment solution yields
substantively distinct, cleanly separated profiles
(Section~\ref{sec-measurement}), whereas the marginal BIC gain at
\(K = 6\) purchases a further split at the cost of entropy. Consistent
with the flatness of the criterion, the unweighted benchmark's own BIC
minimizes at \(K = 6\) as well; enumeration under weighting and without
it point to the same region. We emphasize that \(K\) is an analyst
decision recorded in configuration, not an automated selection, and the
full enumeration evidence is reported so readers can weigh the
neighboring solution.

\section{The measurement model}\label{sec-measurement}

The chosen model converged with weighted log-likelihood \(-29{,}047.7\)
on the sum-to-\(n\) scale, relative entropy 0.907, and mean maximum
posterior probability 0.939: respondents are, on average, assigned to
their modal segment with 94 percent confidence. Table~\ref{tbl-prev}
reports the design-weighted segment prevalences with JKn intervals
alongside the unweighted benchmark, and Figure~\ref{fig-profiles}
displays the full measurement model, the estimated \(\rho_{jck}\) for
every item, category, and segment.

\begin{longtable}[]{@{}
  >{\raggedright\arraybackslash}p{(\columnwidth - 6\tabcolsep) * \real{0.4062}}
  >{\raggedleft\arraybackslash}p{(\columnwidth - 6\tabcolsep) * \real{0.3021}}
  >{\raggedleft\arraybackslash}p{(\columnwidth - 6\tabcolsep) * \real{0.2188}}
  >{\raggedleft\arraybackslash}p{(\columnwidth - 6\tabcolsep) * \real{0.0729}}@{}}

\caption{\label{tbl-prev}Design-weighted segment prevalences with 95
percent JKn intervals (logit scale, t critical value with 126 design
degrees of freedom), the unweighted poLCA benchmark, and the weighting
gap.}

\tabularnewline

\toprule\noalign{}
\begin{minipage}[b]{\linewidth}\raggedright
Segment
\end{minipage} & \begin{minipage}[b]{\linewidth}\raggedleft
Weighted prevalence {[}95\% CI{]}
\end{minipage} & \begin{minipage}[b]{\linewidth}\raggedleft
Unweighted benchmark
\end{minipage} & \begin{minipage}[b]{\linewidth}\raggedleft
Gap
\end{minipage} \\
\midrule\noalign{}
\endhead
\bottomrule\noalign{}
\endlastfoot
S1: Moderate Institutional Trust & 0.331 {[}0.293, 0.372{]} & 0.336 &
-0.005 \\
S2: Generalized Institutional Distrust & 0.107 {[}0.086, 0.132{]} &
0.107 & +0.000 \\
S3: Presidential and Military Trust & 0.242 {[}0.204, 0.285{]} & 0.239 &
+0.004 \\
S4: Very High National Trust & 0.111 {[}0.093, 0.134{]} & 0.111 &
-0.000 \\
S5: Broad Institutional Trust & 0.208 {[}0.173, 0.249{]} & 0.207 &
+0.001 \\

\end{longtable}

\begin{figure}

\centering{

\includegraphics{survey_lca_paper_files/figure-pdf/fig-profiles-1.pdf}

}

\caption{\label{fig-profiles}Conditional response probabilities
P(response category \textbar{} segment) for all thirteen trust items.
Rows are items, columns are the seven response categories from 1 (nada)
to 7 (mucho).}

\end{figure}%

The five profiles order naturally on overall trust with two
configurational departures. Segment S5, Broad Institutional Trust
(weighted prevalence 0.208), concentrates mass on categories 5 and 6
across every institution. Segment S4, Very High National Trust (0.111),
is its extreme counterpart, with probability 0.82 of answering
\emph{mucho} on trust in the president and 0.70 or above on the public
ministry, armed forces, supreme court, and electoral tribunal. Segment
S2, Generalized Institutional Distrust (0.107), mirrors it at the other
pole, with probability 0.85 of answering \emph{nada} on political
parties and 0.76 on the national police; only the president and armed
forces draw any appreciable mass off the floor. The two largest segments
are configurational rather than uniform. Segment S1, Moderate
Institutional Trust (0.331), holds the middle categories on the justice
and audit institutions while shifting upward on the president and armed
forces. Segment S3, Presidential and Military Trust (0.242), sharpens
that asymmetry: substantial trust in the president (probability 0.21 at
\emph{mucho}, mass concentrated at 5 to 7) and the armed forces, against
firm distrust of political parties (0.41 at \emph{nada}) and the police.
The configuration, not the level, separates S1 from S3, and it is
precisely this configurational structure that a summed trust scale would
average away.

The weighted and unweighted prevalences differ by at most 0.005
(Table~\ref{tbl-prev}): the released weights for this single-country
file are nearly constant, so weighting shifts point estimates little
here. The design still matters for inference: the JKn intervals reflect
the clustering, and the interval for S1, for example, spans 0.293 to
0.372, wider than an SRS binomial interval at \(n = 1{,}430\) would
suggest.

\section{Diagnostics}\label{sec-results-diagnostics}

\begin{figure}

\centering{

\includegraphics{survey_lca_paper_files/figure-pdf/fig-discrim-1.pdf}

}

\caption{\label{fig-discrim}Item discrimination: the mean over segment
pairs of the total variation distance between conditional response
distributions.}

\end{figure}%

Every item earns its place. Discrimination (Figure~\ref{fig-discrim})
ranges from 0.452 (president) to 0.754 (justice system); no item
approaches the region where removal would be considered. The BIC screen
of Section~\ref{sec-diagnostics} agrees: \(\Delta\mathrm{BIC}_j\) is
positive for all thirteen items, from \(+320.0\) (president) to
\(+1{,}419.3\) (justice system), so BIC prefers retaining every item's
class dependence, by margins that are decisive throughout. That the
president and armed forces items discriminate least is itself
informative: they are the institutions on which otherwise different
segments agree.

\begin{longtable}[]{@{}llrr@{}}

\caption{\label{tbl-bvr}Largest bivariate residuals. BVR is a
descriptive ranking under the weighted pseudo-likelihood; no chi-square
reference distribution applies.}

\tabularnewline

\toprule\noalign{}
Item A & Item B & Pair df & BVR \\
\midrule\noalign{}
\endhead
\bottomrule\noalign{}
\endlastfoot
public\_ministry & auditor & 36 & 5.75 \\
armed\_forces & legislature & 36 & 5.54 \\
justice\_system & supreme\_court & 36 & 3.21 \\
police & political\_parties & 36 & 2.81 \\
president & supreme\_court & 36 & 2.48 \\

\end{longtable}

The bivariate residuals (Table~\ref{tbl-bvr}) identify two item pairs
standing above the rest: public ministry with auditor (5.75) and armed
forces with legislature (5.54), against a table otherwise at or below
3.21. These are institutionally adjacent pairs, two
audit-and-prosecution bodies and two national-state bodies, whose
residual association the five segments do not fully absorb. Because no
reference distribution is valid here, we read these as a ranked caution
rather than a rejection, one consistent with BIC's mild preference for
the neighboring six-segment solution, and we flag the two pairs as the
natural first check in any sensitivity refit.

\section{Validation and assignment
coverage}\label{sec-results-validation}

At unit weights the pipeline's EM and \texttt{poLCA} agree exactly: the
maximum absolute prevalence gap is 0.00000, the maximum
conditional-probability gap is 0.00000, and the two log-likelihoods are
identical to reported precision (-29046.7917). The implementations
maximize the same function and find the same optimum, which certifies
the weighted estimator's machinery before any weighted result is
interpreted.

Prediction under Equation~\ref{eq-partial} assigned a segment to 1,618
of 1,622 respondents: all 1,430 complete cases plus 188 item-partial
respondents answering at least the seven-item floor. Four respondents
fell below the floor and were left unassigned rather than guessed. The
delivered data product is a labelled SPSS file carrying the numeric
segment with value labels, the label text, and the assignment
diagnostics, together with CSV exports of the full measurement model
with intervals and the demographic composition estimates.

\section{Demographic composition}\label{sec-results-demographics}

BCH-corrected composition was estimated for seven covariates; we
highlight the two with the sharpest structure, presidential job approval
and satisfaction with democracy, and summarize the remainder. Full
tables for all covariates accompany the replication materials.

\begin{longtable}[]{@{}
  >{\raggedright\arraybackslash}p{(\columnwidth - 10\tabcolsep) * \real{0.1858}}
  >{\raggedleft\arraybackslash}p{(\columnwidth - 10\tabcolsep) * \real{0.1593}}
  >{\raggedleft\arraybackslash}p{(\columnwidth - 10\tabcolsep) * \real{0.1593}}
  >{\raggedleft\arraybackslash}p{(\columnwidth - 10\tabcolsep) * \real{0.1593}}
  >{\raggedleft\arraybackslash}p{(\columnwidth - 10\tabcolsep) * \real{0.1770}}
  >{\raggedleft\arraybackslash}p{(\columnwidth - 10\tabcolsep) * \real{0.1593}}@{}}

\caption{\label{tbl-prez}BCH-corrected segment prevalence within
presidential job approval, with 95 percent JKn intervals on the logit
scale. Rows sum to one across segments. One cell (Bad approval, S4) has
a slightly negative BCH point estimate, a documented property of the
correction; its interval is truncated Wald.}

\tabularnewline

\toprule\noalign{}
\begin{minipage}[b]{\linewidth}\raggedright
Level
\end{minipage} & \begin{minipage}[b]{\linewidth}\raggedleft
S1
\end{minipage} & \begin{minipage}[b]{\linewidth}\raggedleft
S2
\end{minipage} & \begin{minipage}[b]{\linewidth}\raggedleft
S3
\end{minipage} & \begin{minipage}[b]{\linewidth}\raggedleft
S4
\end{minipage} & \begin{minipage}[b]{\linewidth}\raggedleft
S5
\end{minipage} \\
\midrule\noalign{}
\endhead
\bottomrule\noalign{}
\endlastfoot
Very bad & 0.21 {[}0.11, 0.38{]} & 0.52 {[}0.35, 0.68{]} & 0.14 {[}0.05,
0.33{]} & 0.05 {[}0.01, 0.33{]} & 0.08 {[}0.02, 0.24{]} \\
Bad & 0.30 {[}0.19, 0.44{]} & 0.31 {[}0.20, 0.44{]} & 0.33 {[}0.21,
0.49{]} & -0.00 {[}0.00, -0.00{]} & 0.06 {[}0.02, 0.17{]} \\
Neither good nor bad & 0.39 {[}0.32, 0.46{]} & 0.16 {[}0.13, 0.20{]} &
0.26 {[}0.21, 0.32{]} & 0.07 {[}0.04, 0.10{]} & 0.12 {[}0.08, 0.18{]} \\
Good & 0.33 {[}0.30, 0.37{]} & 0.06 {[}0.04, 0.09{]} & 0.24 {[}0.21,
0.28{]} & 0.10 {[}0.08, 0.13{]} & 0.26 {[}0.22, 0.30{]} \\
Very good & 0.30 {[}0.25, 0.35{]} & 0.07 {[}0.05, 0.10{]} & 0.21
{[}0.17, 0.26{]} & 0.19 {[}0.16, 0.24{]} & 0.23 {[}0.18, 0.28{]} \\

\end{longtable}

\begin{figure}

\centering{

\includegraphics{survey_lca_paper_files/figure-pdf/fig-prez-1.pdf}

}

\caption{\label{fig-prez}BCH-corrected segment composition by
presidential job approval (95 percent JKn intervals).}

\end{figure}%

Composition is strongly ordered in presidential approval
(Table~\ref{tbl-prez}, Figure~\ref{fig-prez}). Among respondents rating
the president's performance very bad, Generalized Institutional Distrust
reaches 0.52 {[}0.35, 0.68{]}; among those rating it very good, that
segment falls to 0.07 while Very High National Trust rises to 0.19
{[}0.16, 0.24{]}. The correction's behavior at sparse cells is visible
in the Bad approval row, where the S4 estimate is a numerically negative
zero with a truncated interval: no respondents rating the president
badly belong, after correction, to the segment defined by near-uniform
\emph{mucho} responses that include the president item, exactly as
coherence requires. Satisfaction with democracy behaves symmetrically:
Generalized Distrust is 0.34 {[}0.25, 0.45{]} among the very
dissatisfied against 0.06 among the very satisfied, while Very High
National Trust is 0.34 {[}0.25, 0.44{]} among the very satisfied against
0.05 among the very dissatisfied.

The remaining covariates show more modest structure. Presidential and
Military Trust is more prevalent in urban areas (0.26 {[}0.23, 0.29{]})
than rural (0.18 {[}0.14, 0.22{]}) and declines with age, from 0.25
among respondents 18 to 29 to 0.17 among those 60 and over, while Broad
Institutional Trust moves in the opposite direction with age.
Respondents with no formal education are estimated predominantly in Very
High National Trust (0.41 {[}0.22, 0.64{]}), though this cell rests on
few observations and its interval is correspondingly wide. Differences
by sex and employment status are small, with overlapping intervals
throughout.

Discussion

The pipeline demonstrated here closes a specific gap: design-weighted
latent class estimation with fully integrated, design-based replicate
inference, in open R code, validated against an independent
implementation, and carried through to bias-adjusted profiling and a
labelled data product. Each component is individually established,
pseudo-maximum likelihood for LCA under complex sampling
\citep{patterson2002}, weight-normalized information criteria
\citep{lumley2015}, JKn replication \citep{wolter2007}, and the BCH
correction \citep{bolck2004, vermunt2010}; the contribution is their
disciplined integration, with the seams between components, alignment of
replicate fits, a single frozen validation verdict, one interval
construction throughout, treated as first-class design decisions rather
than afterthoughts.

The application illustrates why the machinery matters substantively.
Institutional trust in Mexico is not unidimensional: two of the five
segments, together a third of the population, are defined by the
configuration of trust, elevated for the president and armed forces
against distrust of parties and police, rather than by its level, and
segment composition moves sharply with presidential approval. A summed
trust scale would compress exactly this structure. The near-constancy of
the released weights in this file makes the weighting correction small
for point estimates while leaving the design correction to inference
fully operative; in surveys with more variable weights the same pipeline
applies unchanged with larger weighting gaps.

Several limitations bound the claims. The pseudo-likelihood can be
multimodal with near-tied optima; the reported solution is the best mode
under many seeded starts, and JKn replication expresses within-mode
design variance only. The BCH step conditions on the fitted measurement
model; with relative entropy 0.907 the resulting understatement of
variance is minor \citep{bakk2014}, but it is an understatement.
Estimation uses complete cases, with model-based assignment rather than
imputation extending results to item-partial respondents above the
information floor. The choice of \(K\) is an argued analyst decision
within a flat region of the criterion, and we report the full
enumeration evidence precisely so that readers can weigh the neighboring
six-segment solution. Finally, language-model-drafted labels are a
documented convenience with a certified instrument and full analyst
override, not a measurement claim; all substantive interpretation rests
on the estimated conditional response probabilities.

Replication materials, including all source code, the rendered analysis
documents, and the frozen output objects from which every table and
figure in this article is generated, are available in the project
repository.

References

\renewcommand{\bibsection}{}
\bibliography{bibliography.bib}





\end{document}
